\subsection{Truncamiento por irreducibles y control del resto}
\label{ix-add:corte-primos}

La profundidad de ocupación $K$ y el límite del irreducible $Y$ definen
aproximaciones diferentes de la misma representación aritmética. Las cotas
anteriores conservan todos los irreducibles y truncan la profundidad. Ahora
se mantienen completas las contribuciones locales y se restringen los modos
producidos por el selector a $p\leq Y$, con $Y\geq2$ entero. Se definen
\[
 B_1(Y)=\gamma_0+\sum_{p\leq Y}
       \left(\log(1-p^{-1})+p^{-1}\right),\qquad
 C_2(Y)=\prod_{2<p\leq Y}\left(1-\frac1{(p-1)^2}\right).
\]
El producto vacío vale uno. La parte finita $\gamma_0$ es el mismo objeto
denotado por $\gamma$ en la fórmula de Meissel--Mertens; su evaluación con
resto se desarrolla en la sección~\ref{lectura:manuscrito-04}.

\HMTresultado{Proposición}{Encierros por corte de irreducibles}{ix-add:prime-cut-bound}
Para todo entero $Y\geq2$,
\[
 B_1(Y)-\frac1{2Y}\leq B_1\leq B_1(Y),\qquad
 \frac{Y-1}{Y}C_2(Y)\leq C_2\leq C_2(Y).
\]
\textbf{Demostración.}
Para cada $p\geq2$, la contribución que se resta de $B_1(Y)$ es
\[
 d_p=-\log(1-p^{-1})-p^{-1}
     =\sum_{k\geq2}\frac1{kp^k},\qquad
 0<d_p\leq\frac1{2p(p-1)}.
\]
La última desigualdad utiliza $1/k\leq1/2$ y la suma geométrica.
Por comparación positiva y telescopía,
\[
 0\leq B_1(Y)-B_1=\sum_{p>Y}d_p
 \leq\frac12\sum_{n>Y}\frac1{n(n-1)}=\frac1{2Y}.
\]
Cada factor de la cola de $C_2$ pertenece a $(0,1)$. Incluir todos los
enteros $n>Y$ disminuye el producto. Para $M>Y$,
\[
 \prod_{n=Y+1}^{M}\left(1-\frac1{(n-1)^2}\right)
 =\prod_{n=Y+1}^{M}\frac{n(n-2)}{(n-1)^2}
 =\frac{M(Y-1)}{Y(M-1)}.
\]
El límite es $(Y-1)/Y$. La convergencia del producto primo y la
comparación anterior dan el segundo encierro. \hfill$\square$

Al aumentar $Y$, $B_1(Y)$ y $C_2(Y)$ decrecen hacia sus valores respectivos.
Las anchuras garantizadas son $1/(2Y)$ y $C_2(Y)/Y$. Los encierros por
profundidad y por irreducible pueden intersectarse para refinar la evaluación;
cada uno conserva su parámetro y su demostración. El producto $C_2$
corresponde al sector $p>2$; el factor local del modo dos mantiene la
normalización $2C_2$ de la serie singular completa.
